
\documentclass{article}
\usepackage{colortbl}
\usepackage{makecell}
\usepackage{multirow}
\usepackage{supertabular}

\begin{document}

\newcounter{utterance}

\centering \large Interaction Transcript for game `referencegame', experiment `line\_grids\_columns', episode 1 with NEW\_Reflect\_Agent.
\vspace{24pt}

{ \footnotesize  \setcounter{utterance}{1}
\setlength{\tabcolsep}{0pt}
\begin{supertabular}{c@{$\;$}|p{.15\linewidth}@{}p{.15\linewidth}p{.15\linewidth}p{.15\linewidth}p{.15\linewidth}p{.15\linewidth}}
    \# & $\;$A & \multicolumn{4}{c}{Game Master} & $\;\:$B\\
    \hline

    \theutterance \stepcounter{utterance}  
    & & \multicolumn{4}{p{0.6\linewidth}}{
        \cellcolor[rgb]{0.9,0.9,0.9}{
            \makecell[{{p{\linewidth}}}]{
                \texttt{\tiny{[P1$\langle$GM]}}
                \texttt{You are given three grids, where each of them is 5 by 5 in size.} \\
\texttt{Grids have empty cells marked with "▢" and filled cells marked with "X".} \\
\texttt{Your task is to generate a referring expression that best describes the target grid while distinguishing it from the two other distractor grids.} \\
\texttt{The first grid is the target grid, and the following two grids are the distractors.} \\
\\ 
\texttt{Target grid:} \\
\\ 
\texttt{□ □ X X X} \\
\texttt{□ □ X X X} \\
\texttt{□ □ X X X} \\
\texttt{□ □ X X X} \\
\texttt{□ □ X X X} \\
\\ 
\texttt{Distractor grid 1:} \\
\\ 
\texttt{X □ □ X X} \\
\texttt{X □ □ X X} \\
\texttt{X □ □ X X} \\
\texttt{X □ □ X X} \\
\texttt{X □ □ X X} \\
\\ 
\texttt{Distractor grid 2:} \\
\\ 
\texttt{□ X X X □} \\
\texttt{□ X X X □} \\
\texttt{□ X X X □} \\
\texttt{□ X X X □} \\
\texttt{□ X X X □} \\
\\ 
\texttt{Instruction: Describe the target grid.} \\
\texttt{Generate the referring expression starting with the tag "Expression: " for the given target grid. Omit any other text.} \\
            }
        }
    }
    & \\ \\

    \theutterance \stepcounter{utterance}  
    & \multicolumn{4}{p{0.45\linewidth}}{
        \cellcolor[rgb]{0.95,0.95,0.95}{
            \makecell[{{p{\linewidth}}}]{
                \texttt{\tiny{[P1$|$internal]}}
                \texttt{Reflections from previous episodes:} \\
\texttt{{-}} \\
\texttt{ Failed. I described the target as having "two leftmost columns empty" but distractors also had two leftmost columns empty (though with different patterns). Should have specified that the target's two left columns are entirely empty, while distractor 1 has three left columns empty and distractor 2 has two left columns with alternating X's. Next time, be precise about the exact empty/filled pattern in the left columns.} \\
\texttt{{-}} \\
\texttt{ Failed. I described only the target's filled columns, but that did not distinguish it from Distractor 2, which also has two rightmost columns filled. Next time, I must specifically note the third column is empty while the first two columns are filled — that uniquely contrasts with both distractors.} \\
\texttt{{-}} \\
\texttt{ Failed. My expression didn't distinguish the target because I focused only on the left column of X's, which is also present in distractor 2. I should have specified the exact pattern: a vertical column of X's on the left, then a vertical column of empty cells, then two vertical columns of X's. Next time, I must describe the full unique column pattern, not just one part.} \\
            }
        }
    }
    & & \\ \\

    \theutterance \stepcounter{utterance}  
    & \multicolumn{4}{p{0.45\linewidth}}{
        \cellcolor[rgb]{0.95,0.95,0.95}{
            \makecell[{{p{\linewidth}}}]{
                \texttt{\tiny{[P1$|$internal]}}
                \texttt{{[}attempt 1 – valid{]}} \\
\texttt{Expression: The grid has the first two columns entirely empty, the third column filled with X's, and the last two columns filled with X's.} \\
            }
        }
    }
    & & \\ \\

    \theutterance \stepcounter{utterance}  
    & \multicolumn{4}{p{0.6\linewidth}}{
        \cellcolor[rgb]{0.9,0.9,0.9}{
            \makecell[{{p{\linewidth}}}]{
                \texttt{\tiny{[P1$\rangle$GM]}}
                \texttt{Expression: The grid has the first two columns entirely empty, the third column filled with X's, and the last two columns filled with X's.} \\
            }
        }
    }
    & & \\ \\

    \theutterance \stepcounter{utterance}  
    & & & \multicolumn{2}{p{0.3\linewidth}}{
        \cellcolor[rgb]{0.9,0.9,0.9}{
            \makecell[{{p{\linewidth}}}]{
                \texttt{\tiny{[GM$|$GM]}}
                \texttt{Expression: The grid has the first two columns entirely empty, the third column filled with X's, and the last two columns filled with X's.} \\
            }
        }
    }
    & & \\ \\

    \theutterance \stepcounter{utterance}  
    & & \multicolumn{4}{p{0.6\linewidth}}{
        \cellcolor[rgb]{0.9,0.9,0.9}{
            \makecell[{{p{\linewidth}}}]{
                \texttt{\tiny{[GM$\rangle$P2]}}
                \texttt{You are given three grids, where each of them is 5 by 5 in size.} \\
\texttt{Grids have empty cells marked with "▢" and filled cells marked with "X".} \\
\texttt{You are also given a referring expression that describes one of the given grids.} \\
\texttt{Your task is to select the grid that matches the given referring expression.} \\
\texttt{Generate only the number (in text) of the grid that matches the given expression by selecting first, second, or third.} \\
\\ 
\texttt{First:} \\
\\ 
\texttt{□ □ X X X} \\
\texttt{□ □ X X X} \\
\texttt{□ □ X X X} \\
\texttt{□ □ X X X} \\
\texttt{□ □ X X X} \\
\\ 
\texttt{Second:} \\
\\ 
\texttt{X □ □ X X} \\
\texttt{X □ □ X X} \\
\texttt{X □ □ X X} \\
\texttt{X □ □ X X} \\
\texttt{X □ □ X X} \\
\\ 
\texttt{Third:} \\
\\ 
\texttt{□ X X X □} \\
\texttt{□ X X X □} \\
\texttt{□ X X X □} \\
\texttt{□ X X X □} \\
\texttt{□ X X X □} \\
\\ 
\texttt{Expression: The grid has the first two columns entirely empty, the third column filled with X's, and the last two columns filled with X's.} \\
\texttt{Question: Which grid does the expression refer to?} \\
\texttt{Start with the tag "Answer: ", followed by your selection. Omit any other text.} \\
            }
        }
    }
    & \\ \\

    \theutterance \stepcounter{utterance}  
    & & & \multicolumn{4}{p{0.45\linewidth}}{
        \cellcolor[rgb]{0.95,0.95,0.95}{
            \makecell[{{p{\linewidth}}}]{
                \texttt{\tiny{[P2$|$internal]}}
                \texttt{Reflections from previous episodes:} \\
\texttt{{-}} \\
\texttt{ Failed. I misinterpreted the expression. It described two full rows of X's at the top and bottom, with all other rows empty. None of the grids match that description. I should have recognized the mismatch and possibly concluded "none" or re{-}checked. Next time, I will carefully visualize the described pattern and compare it literally to each grid before answering.} \\
\texttt{{-}} \\
\texttt{ Failed. I misinterpreted the expression, focusing on "two full rows of X's" as completely filled rows, but the description actually specified "no other X's in those rows" meaning only those rows have X's in all columns. I should have checked each grid's row patterns more carefully against the exact column conditions given. Next time, I will parse the expression step{-}by{-}step, verifying each row and column constraint directly on the grids.} \\
\texttt{{-}} \\
\texttt{ Failed. The expression was overly complex and repetitive, causing confusion instead of clarity. I tried to parse it literally but got lost in the contradictions and circular descriptions. Next time, I would ignore the convoluted text and simply match the described pattern directly to the visual grids, which would have shown that none of the grids fully matched the described structure.} \\
            }
        }
    }
     \\ \\

    \theutterance \stepcounter{utterance}  
    & & & \multicolumn{4}{p{0.45\linewidth}}{
        \cellcolor[rgb]{0.95,0.95,0.95}{
            \makecell[{{p{\linewidth}}}]{
                \texttt{\tiny{[P2$|$internal]}}
                \texttt{{[}attempt 1 – valid{]}} \\
\texttt{Answer: third} \\
            }
        }
    }
     \\ \\

    \theutterance \stepcounter{utterance}  
    & & & \multicolumn{4}{p{0.6\linewidth}}{
        \cellcolor[rgb]{0.9,0.9,0.9}{
            \makecell[{{p{\linewidth}}}]{
                \texttt{\tiny{[GM$\langle$P2]}}
                \texttt{Answer: third} \\
            }
        }
    }
     \\ \\

    \theutterance \stepcounter{utterance}  
    & & & \multicolumn{2}{p{0.3\linewidth}}{
        \cellcolor[rgb]{0.9,0.9,0.9}{
            \makecell[{{p{\linewidth}}}]{
                \texttt{\tiny{[GM$|$GM]}}
                \texttt{third} \\
            }
        }
    }
    & & \\ \\

\end{supertabular}
}

\end{document}
